\section{Liste Concatenate: Concetti Base (Lab 12)}

\begin{concept}{Stack (LIFO) vs Queue (FIFO)}
\textbf{Stack:} Pila di piatti. Ultimo dentro, primo fuori (LIFO).
\textbf{Queue:} Fila alle poste. Primo dentro, primo fuori (FIFO).
\end{concept}

\begin{center}
\begin{tikzpicture}[
    stackbox/.style={draw, thick, fill=niceBlue!10, minimum width=2.5cm, minimum height=0.8cm, font=\ttfamily},
    arrow/.style={-{Stealth[length=3mm]}, thick, color=codeBlue}
]
% Contenitore Stack
\draw[thick] (-1.5, 0) -- (-1.5, -3.5) -- (1.5, -3.5) -- (1.5, 0);
% Elementi
\node[stackbox] (item1) at (0, -3.0) {index 0: 10};
\node[stackbox] (item2) at (0, -2.1) {index 1: 20};
\node[stackbox, fill=jsYellow!20] (item3) at (0, -1.2) {index 2: 30};
% Frecce
\draw[arrow] (2.5, -0.5) to[out=180, in=0] node[midway, above, font=\small\bfseries] {PUSH} (item3.east);
\end{tikzpicture}
\end{center}

\subsection{Cosa sono le Liste Concatenate?}
A differenza degli Array, dove gli elementi sono vicini in memoria e usiamo gli indici (`a[0]`), nelle Liste gli elementi sono sparsi.
Ogni elemento (**Nodo**) contiene:
\begin{itemize}
    \item **val**: Il dato (numero, stringa, etc.)
    \item **next**: Un riferimento (freccia) al prossimo nodo. L'ultimo punta a `null`.
\end{itemize}

Struttura: `Testa -> [5 | ->] -> [3 | ->] -> [6 | null]`

\subsubsection{Immaginiamolo visivamente}
\begin{verbatim}
       HEAD
        |
        v
    +------+      +------+      +------+
    | 5    | ---> | 3    | ---> | 6    | ---> null
    +------+      +------+      +------+
\end{verbatim}

\subsection{Definizione in JavaScript}
\begin{concept}{Struttura Nodo}
Non serve una classe speciale, bastano gli oggetti!
Ogni nodo ha un valore e un puntatore al prossimo.
\end{concept}
\begin{jsbox}
let nodo3 = { val: 6, next: null };
let nodo2 = { val: 3, next: nodo3 };
let head  = { val: 5, next: nodo2 };
\end{jsbox}

\subsection{Operazioni Base (Ricorsive)}

\subsubsection{1. Trovare un elemento (listFind)}
\begin{concept}{Scansione Lineare}
Scorriamo la lista finché non troviamo il valore o finiamo (null).
\end{concept}
\begin{jsbox}
function listFind(nodo, valore) {
    if (!nodo) return null; // Lista finita, non trovato
    if (nodo.val === valore) return nodo; // Trovato!
    
    // Cerca nel resto della lista
    return listFind(nodo.next, valore);
}
\end{jsbox}

\subsubsection{2. Aggiungere in fondo (listPush)}
\begin{concept}{Inserimento in Coda}
Dobbiamo "scendere" fino all'ultimo nodo (quello con next: null) e attaccarne uno nuovo.
\end{concept}
\begin{jsbox}
function listPush(nodo, valore) {
    // Caso base: se il nodo è null, creo il nuovo nodo qui
    if (!nodo) {
        return { val: valore, next: null };
    }
    
    // Passo ricorsivo: ricostruisco la lista attaccando il risultato
    nodo.next = listPush(nodo.next, valore);
    return nodo;
}
\end{jsbox}

\subsubsection{3. Rimuovere l'ultimo (listPop)}
\begin{concept}{Rimozione Ultimo}
Restituisce una coppia: `[ListaAggiornata, ValoreRimosso]`.
\end{concept}
\begin{jsbox}
function listPop(nodo) {
    if (!nodo) return [null, undefined]; // Lista vuota
    
    // Se sono l'ultimo nodo (il mio next è null)
    if (!nodo.next) {
        return [null, nodo.val]; // Mi rimuovo (ritorno null al padre)
    }
    
    // Altrimenti vado avanti
    let [nuovaCoda, valRimosso] = listPop(nodo.next);
    nodo.next = nuovaCoda; // Aggiorno il mio puntatore
    return [nodo, valRimosso];
}
\end{jsbox}
